\section{Rangkuman}
\begin{itemize}
    \item Serangan adalah setiap usaha pihak lawan untuk menemukan kunci atau plainteks; Prinsip Kerckhoffs menegaskan algoritma publik dan hanya kunci yang rahasia.
    \item Berdasarkan informasi yang dimiliki penyerang, serangan dibedakan menjadi ciphertext-only, known-plaintext, chosen-plaintext, adaptive chosen-plaintext, dan chosen-ciphertext. Urutannya menunjukkan kekuatan penyerang yang meningkat.
    \item Ciphertext-only adalah yang tersulit karena penyerang hanya berbekal cipherteks, misalnya cipherteks Caesar \code{KHOORZRUOG} yang dipecahkan dengan mencoba 25 pergeseran.
    \item Pada chosen-plaintext, penyerang memilih sendiri plainteksnya; pada Vigen\`ere, mengirim \code{AAAAA} membuat cipherteks yang keluar sama dengan kuncinya sendiri. Pada chosen-ciphertext, penyerang memakai sistem dekripsi sebagai \textit{oracle}.
    \item Berdasarkan tujuan akhirnya, serangan meliputi key-recovery, plaintext-recovery, distinguishing, forgery, dan replay. Keberhasilan pada satu tujuan belum tentu berarti keberhasilan pada tujuan lain.
    \item Serangan pasif hanya menyadap tanpa mengubah data (intersepsi, analisis lalu lintas), sedangkan serangan aktif mengintervensi komunikasi (interupsi, fabrikasi, modifikasi, pemutaran ulang, MITM).
    \item Serangan pasif sulit dideteksi karena tidak meninggalkan jejak perubahan apa pun; serangan aktif dapat dideteksi justru karena mengubah keadaan sistem.
    \item Alat seperti Wireshark memperlihatkan bahwa pada lalu lintas yang tidak terenkripsi, nama pengguna dan kata sandi terbaca langsung tanpa perlu dipecahkan.
    \item Selain penyadapan perangkat lunak, dikenal penyadapan kabel, penyadapan elektromagnetik (radiasi perangkat, jarak kurang dari dua meter), dan penyadapan akustik melalui vibrometri.
    \item Serangan MITM berbahaya karena kedua pihak sama-sama merasa berkomunikasi langsung dengan pihak yang benar; penangkalnya adalah otentikasi, bukan sekadar enkripsi.
    \item Sistem disebut aman secara komputasi bila biaya atau waktu pemecahannya melampaui nilai dan masa berlaku informasi. Keamanan komputasi berbeda dari keamanan tanpa syarat, yang hanya dipenuhi oleh one-time pad.
    \item Serangan saluran samping mengeksploitasi kebocoran fisik perangkat, bukan kelemahan matematis algoritmanya: waktu eksekusi, konsumsi daya (SPA/DPA), radiasi elektromagnetik, dan penyuntikan kesalahan (DFA).
    \item Penangkalan serangan saluran samping pada dasarnya adalah membuat perilaku fisik perangkat tidak lagi bergantung pada nilai rahasia yang sedang diolahnya.
\end{itemize}
